\documentclass[10pt,a4paper]{amsart}
\usepackage{amsmath,amssymb,amsthm,graphicx}
\usepackage{verbatim}
\newenvironment{reference}{\comment}{\endcomment}
%\newenvironment{reference}{}{}
\newenvironment{slogan}{\comment}{\endcomment}
\newenvironment{history}{\comment}{\endcomment}

% For commutative diagrams we use Xy-pic
\usepackage[all]{xy}

% We use 2cell for 2-commutative diagrams.
\xyoption{2cell}
% Theorem environments.
%
\theoremstyle{plain}
\newtheorem{theorem}[subsection]{Theorem}
\newtheorem{proposition}[subsection]{Proposition}
\newtheorem{lemma}[subsection]{Lemma}

\theoremstyle{definition}
\newtheorem{definition}[subsection]{Definition}
\newtheorem{example}[subsection]{Example}
\newtheorem{exercise}[subsection]{Exercise}
\newtheorem{situation}[subsection]{Situation}

\theoremstyle{remark}
\newtheorem{remark}[subsection]{Remark}
\newtheorem{remarks}[subsection]{Remarks}

\numberwithin{equation}{subsection}

% Macros
%
\def\lim{\mathop{\mathrm{lim}}\nolimits}
\def\colim{\mathop{\mathrm{colim}}\nolimits}
\def\Spec{\mathop{\mathrm{Spec}}}
\def\Hom{\mathop{\mathrm{Hom}}\nolimits}
\def\Ext{\mathop{\mathrm{Ext}}\nolimits}
\def\SheafHom{\mathop{\mathcal{H}\!\mathit{om}}\nolimits}
\def\SheafExt{\mathop{\mathcal{E}\!\mathit{xt}}\nolimits}
\def\Sch{\mathit{Sch}}
\def\Mor{\mathop{\mathrm{Mor}}\nolimits}
\def\Ob{\mathop{\mathrm{Ob}}\nolimits}
\def\Sh{\mathop{\mathit{Sh}}\nolimits}
\def\NL{\mathop{N\!L}\nolimits}
\def\CH{\mathop{\mathrm{CH}}\nolimits}
\def\proetale{{pro\text{-}\acute{e}tale}}
\def\etale{{\acute{e}tale}}
\def\QCoh{\mathit{QCoh}}
\def\Ker{\mathop{\mathrm{Ker}}}
\def\Im{\mathop{\mathrm{Im}}}
\def\Coker{\mathop{\mathrm{Coker}}}
\def\Coim{\mathop{\mathrm{Coim}}}

% Boxtimes
%
\DeclareMathSymbol{\boxtimes}{\mathbin}{AMSa}{"02}

%
% Macros for moduli stacks/spaces
%
\def\QCohstack{\mathcal{QC}\!\mathit{oh}}
\def\Cohstack{\mathcal{C}\!\mathit{oh}}
\def\Spacesstack{\mathcal{S}\!\mathit{paces}}
\def\Quotfunctor{\mathrm{Quot}}
\def\Hilbfunctor{\mathrm{Hilb}}
\def\Curvesstack{\mathcal{C}\!\mathit{urves}}
\def\Polarizedstack{\mathcal{P}\!\mathit{olarized}}
\def\Complexesstack{\mathcal{C}\!\mathit{omplexes}}
% \Pic is the operator that assigns to X its picard group, usage \Pic(X)
% \Picardstack_{X/B} denotes the Picard stack of X over B
% \Picardfunctor_{X/B} denotes the Picard functor of X over B
\def\Pic{\mathop{\mathrm{Pic}}\nolimits}
\def\Picardstack{\mathcal{P}\!\mathit{ic}}
\def\Picardfunctor{\mathrm{Pic}}
\def\Deformationcategory{\mathcal{D}\!\mathit{ef}}
\usepackage{cleveref}
\crefname{originalreference}{original reference}{original references}
\Crefname{originalreference}{Original reference}{Original references}
\makeatletter
\newlabel{lemma-H1-injective}{{0AXB}{source}{}{sourcepointer-1}{}}
\newlabel{lemma-H1-injective@cref}{{[originalreference][0][]0AXB}{[1][0][]source}{}{}{}}
\makeatother
\crefname{originallemma}{original lemma}{original lemmas}
\Crefname{originallemma}{Original lemma}{Original lemmas}
\makeatletter
\newlabel{more-algebra-lemma-divides-radical}{{0AXI}{source}{}{sourcepointer-2}{}}
\newlabel{more-algebra-lemma-divides-radical@cref}{{[originallemma][0][]0AXI}{[1][0][]source}{}{}{}}
\makeatother
\crefname{originallemma}{original lemma}{original lemmas}
\Crefname{originallemma}{Original lemma}{Original lemmas}
\makeatletter
\newlabel{varieties-lemma-finite-in-codim-1}{{0AB7}{source}{}{sourcepointer-3}{}}
\newlabel{varieties-lemma-finite-in-codim-1@cref}{{[originallemma][0][]0AB7}{[1][0][]source}{}{}{}}
\makeatother
\title{Uniform boundedness of the torsion of the first cohomology along modifications of a normal surface singularity}
\author{The Stacks Project Authors}
\date{stacks-project:resolve.tex:local-snapshot-20261012; GNU FDL 1.2 or later}
\begin{document}
\maketitle
\noindent\textit{Editorial note.} All delivered mathematical text is the author's own, copied from the pinned Stacks Project chapter file resolve.tex; no statement or proof step was rewritten, reordered or removed.\par\medskip
\noindent\textit{Editorial note.} The counted proof is the author's complete argument: the identification of the cokernel with the torsion of the first cohomology and the reduction to an element with reduced quotient, the comparison with the normalization of the quotient and with the scheme theoretic closure of the divisor, and the use of the injectivity of the restriction map on the first cohomology to obtain a uniform bound.\par\medskip
\noindent\textit{Editorial note.} The label lemma-bound-a-torsion names the counted statement, and the proof environment immediately following it in the source is delivered verbatim as its proof.\par\medskip
\noindent\textit{Editorial note.} The counted statement is carried by the author's own lemma environment in the chapter source.\par\medskip
\noindent\textit{Editorial note.} Pre-open gate: python operations/agent-authoring/tools/precard\_gate.py --source resolve.tex --label lemma-bound-a-torsion --start 1946 --end 2005 --local-id 28085 returned RESULT=OK (L28085 free, occupancy FREE, public interval clear). Pre-assembly tag self-check: canonical tag resolved and not present in reference/public-items.jsonl.\par\medskip
\noindent\textit{Editorial note.} The statement and proof are copied from the pinned chapter file; the cited inputs are kept as external pointers with their canonical Stacks tags.\par\medskip
\noindent\textit{Allowed standard prerequisites (editorial).}\begin{itemize}
\item Identification of the cokernel with the torsion of the first cohomology. The long exact cohomology sequence of the short exact sequence multiplying by a nonzero element on the structure sheaf. Supplies the identification of the cokernel of the restriction map from the base ring to the sections of the divisor with the elements of the first cohomology annihilated by that element.
\item Reduction to an element generating a radical quotient. More on Algebra, the lemma allowing a nonzero element to be replaced by one whose quotient is reduced and whose powers are divisible by the given element. Supplies the reduction to the case where the quotient by the chosen element is a reduced ring.
\item Finiteness of the normalization of the quotient for Nagata rings, and finiteness in codimension one. Morphisms, the lemma that the normalization of a Nagata scheme is finite, and Varieties, the lemma that a modification of a normal surface is finite in codimension one. Supply the finiteness of the cokernel of the inclusion of the quotient into its normalization and the finiteness of the scheme theoretic closure of the divisor away from the exceptional locus.
\item Injectivity of the restriction of the first cohomology. Same chapter, the lemma asserting the injectivity of the restriction map on the first cohomology of the structure sheaf to the inverse image of the open. Supplies the step showing that a section vanishing on the punctured locus comes from the base ring.
\end{itemize}\medskip
\ifdefined\hypertarget\hypertarget{sourcepointer-1}{}\fi
\noindent\textit{Source reference.} Original reference 0AXB: sources/stacks-resolve/pinned-original-source.tex, source lines 1692--1692. Prerequisite pointer invoked at proof line(s) 1998; it names a standard fact of the same chapter that the counted argument uses, and every step of the counted proof built on it is retained verbatim in the delivered body.\par\medskip
\ifdefined\hypertarget\hypertarget{sourcepointer-2}{}\fi
\noindent\textit{Source reference.} Original lemma 0AXI: sources/stacks-chow/tags.txt, source lines 13756--13756. Prerequisite pointer invoked at proof line(s) 1978; it names a standard fact of a sibling chapter that the counted argument uses, and every step of the counted proof built on it is retained verbatim in the delivered body.\par\medskip
\ifdefined\hypertarget\hypertarget{sourcepointer-3}{}\fi
\noindent\textit{Source reference.} Original lemma 0AB7: sources/stacks-chow/tags.txt, source lines 13010--13010. Prerequisite pointer invoked at proof line(s) 1988; it names a standard fact of a sibling chapter that the counted argument uses, and every step of the counted proof built on it is retained verbatim in the delivered body.\par\medskip
\begin{lemma}
\label{lemma-bound-a-torsion}
Let $(A, \mathfrak m, \kappa)$ be a local normal Nagata domain
of dimension $2$. Let $a \in A$ be nonzero. There exists an integer $N$ such
that for every modification $f : X \to \Spec(A)$ with $X$ normal the
$A$-module
$$
M_{X, a} = \Coker(A \longrightarrow H^0(Z, \mathcal{O}_Z))
$$
where $Z \subset X$ is cut out by $a$ has length bounded by $N$.
\end{lemma}

\begin{proof}
By the short exact sequence
$
0 \to \mathcal{O}_X \xrightarrow{a} \mathcal{O}_X \to \mathcal{O}_Z \to 0
$
we see that
\begin{equation}
\label{equation-a-torsion}
M_{X, a} = H^1(X, \mathcal{O}_X)[a]
\end{equation}
Here $N[a] = \{n \in N \mid an = 0\}$ for an $A$-module $N$. Thus
if $a$ divides $b$, then $M_{X, a} \subset M_{X, b}$.
Suppose that for some $c \in A$ the modules $M_{X, c}$
have bounded length. Then for every $X$ we have an exact sequence
$$
0 \to M_{X, c} \to M_{X, c^2} \to M_{X, c}
$$
where the second arrow is given by multiplication by $c$. Hence we see that
$M_{X, c^2}$ has bounded length as well. Thus it suffices to find a $c \in A$
for which the lemma is true such that $a$ divides $c^n$ for some $n > 0$.
By More on Algebra, Lemma \ref{more-algebra-lemma-divides-radical}
we may assume $A/(a)$ is a reduced ring.

\medskip\noindent
Assume that $A/(a)$ is reduced. Let $A/(a) \subset B$ be the normalization
of $A/(a)$ in its quotient ring. Because $A$ is Nagata, we see that
$\Coker(A \to B)$ is finite. We claim the length of this finite
module is a bound. To see this, consider $f : X \to \Spec(A)$ as in the lemma
and let $Z' \subset Z$ be the scheme theoretic closure of $Z \cap f^{-1}(U)$.
Then $Z' \to \Spec(A/(a))$ is finite for example by Varieties, Lemma
\ref{varieties-lemma-finite-in-codim-1}.
Hence $Z' = \Spec(B')$ with $A/(a) \subset B' \subset B$.
On the other hand, we claim the map
$$
H^0(Z, \mathcal{O}_Z) \to H^0(Z', \mathcal{O}_{Z'})
$$
is injective. Namely, if $s \in H^0(Z, \mathcal{O}_Z)$
is in the kernel, then
the restriction of $s$ to $f^{-1}(U) \cap Z$ is zero.
Hence the image of $s$ in $H^1(X, \mathcal{O}_X)$ vanishes in
$H^1(f^{-1}(U), \mathcal{O}_X)$. By Lemma \ref{lemma-H1-injective}
we see that $s$ comes from an element $\tilde s$ of $A$. But by
assumption $\tilde s$ maps to zero in $B'$ which implies that $s = 0$.
Putting everything together we see that
$M_{X, a}$ is a subquotient of $B'/A$, namely not every element
of $B'$ extends to a global section of $\mathcal{O}_Z$, but in
any case the length of $M_{X, a}$ is bounded by the length of $B/A$.
\end{proof}

\end{document}
